% Jede Aufgabe erhält eine eigene Arbeitsseite; Zusatzraum bei 6 und 8.
\Exercise{1}{Wege zeichnen und die Eltern finden}
\begin{center}\PascalAufgabeStart\end{center}
\textbf{(a)} Zeichne einen Weg nach $(4,2)$ ein und notiere seine
Schrittfolge. Liste dann \emph{alle} solchen Folgen auf. Warum ist
deine Liste vollständig? Markiere jeweils die Positionen der $R$.

\textbf{(b)} Ordne die Folgen nach ihrem letzten Schritt. Von welchen
beiden Stellen kommen sie? Erkläre so $\binom42=\binom31+\binom32$.

\textbf{(c)} Übertrage die Erklärung auf $(n+1,k)$ mit $1\le k\le n+1$.
Warum heißen die Eltern $(n,k-1)$ und $(n,k)$? Was passiert bei $k=0$?

\textbf{(d)} Welche Wege zählen $\binom22,\binom32,\binom42$?
Erkläre den festen Index $2$ und die schräge Lage dieser Stellen im Dreieck.
\ExerciseSpace{40}

\Exercise{2}{Ränder, Nullen und Spiegelung}
\textbf{(a)} Bestimme die folgenden Werte und begründe jeden mit Wegen:
\[
 \binom00,\quad\binom07,\quad\binom90,\quad\binom91,
 \quad\binom98,\quad\binom99,\quad\binom9{10}.
\]

\textbf{(b)} Vertausche in jedem Weg $L$ und $R$. Wohin führt der
Spiegelweg eines Weges nach $(n,k)$? Begründe die Symmetrie durch
Hin- und Rückzuordnung, für $0\le k\le n$.

\textbf{(c)} Warum darf man $\binom35$ nicht mit
$\binom3{3-5}$ gleichsetzen? Erkläre den richtigen Wert im Wegemodell.
\ExerciseSpace{60}

\Exercise{3}{Rechtsschritte auswählen und zählen}
\textbf{(a)} Wie viele Wege nach $(12,3)$ gibt es? Erkläre zunächst,
warum ein solcher Weg durch drei ausgewählte Schrittpositionen
eindeutig bestimmt ist. Berechne die Anzahl dann mit der Fakultätenformel;
kürze vor dem Ausmultiplizieren.

\textbf{(b)} Wie viele Wege führen nach $(12,9)$? Nutze die Spiegelung.

\textbf{(c)} Beweise für jedes $n\ge2$ durch Zählen der beiden
Rechtsschrittpositionen:
\[
 2\binom n2=n(n-1).
\]
Warum zählt die rechte Seite jeden Weg genau zweimal?
Bestätige die Gleichung anschließend durch Kürzen in der Fakultätenformel.
\ExerciseSpace{55}

\Exercise{4}{Wege beim Ausmultiplizieren}
In den fünf Faktoren von $(a+b)^5$ bedeutet $R$: Wähle $a$;
$L$: Wähle $b$.

\textbf{(a)} Welchen Term liefert die Folge $RLLRL$? Wie viele
Schrittfolgen enthalten dieselbe Anzahl $R$ und $L$? Begründe damit
den Koeffizienten dieses Terms.

\textbf{(b)} Schreibe $(a+b)^5$ vollständig aus, geordnet nach
aufsteigenden Potenzen von $a$. Berechne damit $11^5$ für $a=1$, $b=10$.

\textbf{(c)} Erkläre den binomischen Lehrsatz bei $n=0$ durch den
leeren Weg und das leere Produkt. Prüfe dabei auch $a=b=0$ mit der
vereinbarten Potenzkonvention.
\ExerciseSpace{55}

\Exercise{5}{Endpunkte und gefärbte Schritte}
\textbf{(a)} Wie viele Wege mit acht Schritten gibt es insgesamt?
Ordne sie nach ihrem Endpunkt und erkläre so die Summe der Zeile $8$.

\textbf{(b)} Wie viele dieser Wege enden an einer inneren Stelle,
also weder ganz links noch ganz rechts?

\textbf{(c)} Jeder Rechtsschritt erhält nun unabhängig eine von zwei
Farben; Linksschritte bleiben ungefärbt. Zwei Wege gelten als verschieden,
wenn ihre Schrittfolge oder mindestens eine Farbe verschieden ist.
Zähle auf zwei Arten und beweise
\[
 \sum_{k=0}^n\binom nk2^k=3^n.
\]
Wie viele Färbungen besitzt ein Weg nach $(n,k)$? Wie viele Möglichkeiten
gibt es bei jedem einzelnen Schritt? Prüfe $n=0$ und bestätige das
Ergebnis mit dem binomischen Lehrsatz.
\ExerciseSpace{55}

\Exercise{6}{Nach dem letzten Rechtsschritt ordnen}
Wir betrachten alle Wege nach $(n+3,3)$, also mit drei $R$ und $n$ $L$.
Zu zeigen ist
\[
 \sum_{i=0}^n\binom{i+2}2=\binom{n+3}3.
\]
\textbf{(a)} Sei $i$ die Anzahl der $L$ \emph{vor dem letzten $R$}.
Wo endet der Wegteil vor diesem $R$? Wie viele Schritte $L$ folgen
noch? Begründe die Formel, indem du die Wege nach $i$ einteilst.
Zeige auch, wie aus jedem Anfangsweg der vollständige Weg entsteht.

\textbf{(b)} Beweise dieselbe Formel zusätzlich durch Induktion über $n$,
mit Pascal-Rekursion und Randwerten. Zitiere die Diagonalensummenformel
nicht als fertigen Satz.

\textbf{(c)} Berechne damit
\[
 1+3+6+10+15+21+28+36+45+55.
\]
\ExerciseSpace{45}
\ExtraWork{6}{Vor dem letzten $R$ liegen genau zwei $R$ und $i$ $L$;
danach sind nur noch $L$ möglich. Für die Induktion: Trenne den neu
hinzukommenden Summanden ab und suche zwei Pascal-Eltern.}

\Exercise{7}{Wege an einer Zwischenzeile trennen}
\begin{center}\PascalAufgabeSchnitt\end{center}
\textbf{(a)} Zähle für A, B und C getrennt die Hinwege von der Spitze
und die Weiterwege zum Ziel $(5,3)$. Wie viele ganze Wege gehen
jeweils über diese Station? Addiere und vergleiche mit $\binom53$.

\textbf{(b)} Erkläre, warum du innerhalb jeder Station multiplizierst
und über die Stationen addierst. Welche Rolle spielen $m=2$, $n=3$
und $r=3$ in der Vandermonde-Identität? Was zählt $\binom3{3-k}$?

\textbf{(c)} Das neue Ziel sei $(5,6)$. Erkläre mit Wegen und zusätzlich
an \emph{jedem} Summanden, warum
\[
 \sum_{k=0}^6\binom2k\binom3{6-k}=0.
\]
\ExerciseSpace{35}

\Exercise{8}{In der Mitte trennen und spiegeln}
\textbf{(a)} Trenne Wege nach $(8,4)$ nach vier Schritten.
Warum gibt es über jede Station $(4,k)$ genau $\binom4k^2$ Wege?
Berechne die Summe direkt aus Zeile $4$ und durch Zählen aller Wege zum Ziel.

\textbf{(b)} Berechne die Summe benachbarter Produkte direkt:
\[
 \sum_{k=0}^3\binom3k\binom3{k+1}.
\]
Der Faktor $\binom34$ ist $0$. Vergleiche dein Ergebnis mit $\binom64$.

\textbf{(c)} Beweise für jedes $n\in\N$:
\[
 \boxed{\displaystyle
 \sum_{k=0}^n\binom nk\binom n{k+1}=\binom{2n}{n+1}.}
\]
Trenne Wege nach $(2n,n+1)$ nach $n$ Schritten. Wähle diesmal $k$
als Anzahl der \emph{Linksschritte} im ersten Teil. Bestimme die
Rechtsschritte in beiden Teilen und nutze die Spiegelung. Erkläre
Produkt und Summe und prüfe $n=0$.
\ExerciseSpace{50}
\ExtraWork{8}{Im ersten Teil liegen $k$ Linksschritte, also $n-k$
Rechtsschritte. Wie viele der insgesamt $n+1$ Rechtsschritte fehlen
noch? Denke auch an den Fall $k=n$: Kann der zweite Teil dann gelingen?}
